\documentclass[12pt]{ctexbeamer}	%声明文档类型为beamer幻灯片
\mode <presentation>
\usetheme{Madrid}%%Warsaw [secheader]
%\usecolortheme{beaver}
\beamersetaveragebackground{black!10}

\setbeamercovered{transparent}
%\setbeamertemplate{navigation symbols}{}

\usefonttheme{professionalfonts}
\useinnertheme{circles}%{rectangles}
\setbeamertemplate{itemize item}{$\circledast$}%{\checkmark}


\setbeamertemplate{title page}
{
%  \vbox{}
  \vfill
  \begin{centering}
  \includegraphics[height=2.0cm]{ustc_logo_fig_new.eps}
        \vskip0.6em\par%
    \begin{beamercolorbox}[sep=8pt,center,shadow=true,rounded=true]{title}
      \usebeamerfont{title}\inserttitle\par%
      \ifx\insertsubtitle\@empty%
      \else%
        \vskip0.25em%
        {\usebeamerfont{subtitle}\usebeamercolor[fg]{subtitle}\insertsubtitle\par}%
      \fi%
    \end{beamercolorbox}%
    \vskip1em\par
    \begin{beamercolorbox}[sep=8pt,center]{author}
      \usebeamerfont{author}
 %     \insertauthor
      {
        \begin{tabular}{cc}
  答辩人：&\insertauthor\\
  导\quad 师：& \advisor
  \end{tabular}
      }
    \end{beamercolorbox}
    \begin{beamercolorbox}[sep=8pt,center]{institute}
      \usebeamerfont{institute}\insertinstitute
    \end{beamercolorbox}
    \begin{beamercolorbox}[sep=8pt,center]{date}
      \usebeamerfont{date}\insertdate
    \end{beamercolorbox}\vskip0.5em
  %  {\usebeamercolor[fg]{titlegraphic}\inserttitlegraphic\par}
  \end{centering}
  \vfill
}

\hypersetup{pdfpagemode=FullScreen} % makes your presentation go automatically to full screen

\usepackage{fontspec,xunicode,xltxtra}

\usepackage{multimedia}	%让文档支持多媒体
\usepackage{graphics}	%让文档支持图片
\usepackage{hyperref}	%让文档支持超链接
\usepackage{booktabs}	%让文档支持三线表格
\usepackage{amsmath}	%ams可以让文档支持数学公式
\usepackage{amsfonts}
\usepackage{amssymb}
\usepackage{color}
\usepackage{graphicx,psfrag}
\usepackage{epsfig}

\title{话者识别中失配信息补偿理论和方法研究}	%幻灯片标题
\author[刘青松]{刘青松}
\def\advisor{戴蓓蒨教授}

\institute[中国科学技术大学]{
语音信号与信息处理实验室\\
中国科学技术大学，电子科学与技术系
}	%作者单位
\date{\today}

\setcounter{tocdepth}{1}

\begin{document}
\begin{frame}
\titlepage
\end{frame}

\AtBeginSection[]{
  \frame{
    \frametitle{大纲}
    \tableofcontents[currentsection]
  }
}

\begin{frame}
\frametitle{大纲}
\tableofcontents
\end{frame}

%%%%%%%%%%%%%%%%%%%%%%%%%%%%%%%%%%%%%%%%%%%%%%%%%%%%%%%%%%%%%%%%%%%%%%%%%%%%%%%

\section{概述}
\begin{frame}
\frametitle{概述}
研究意义
\end{frame}

%%%%%%%%%%%%%%%%%%%%%%%%%%%%%%%%%%%%%%%%%%%%%%%%%%%%%%%%%%%%%%%%%%%%%%%%%%%%%%%

\section{存在的问题}

\begin{frame}
\frametitle{一二三四五问题}
\begin{itemize}
\item 一二三四五
\item \alert{六七八九十}
\end{itemize}
\pause
一二三四五

\end{frame}

%%%%%%%%%%%%%%%%%%%%%%%%%%%%%%%%%%%%%%%%%%%%%%%%%%%%%%%%%%%%%%%%%%%%%%%%%%%%%%%

\section{方法一}

\begin{frame}
\frametitle{USTC}

\begin{itemize}
\item<1-> USTC,ABCD EFG HIJK LMN OPQ RST UVW XYZ
\item<2-> USTC,abcd efg hijk lmn opq rst uvw xyz
\end{itemize}


\end{frame}

%%%%%%%%%%%%%%%%%%%%%%%%%%%%%%%%%%%%%%%%%%%%%%%%%%%%%%%%%%%%%%%%%%%%%%%%%%%%%%%

\section{方法二}

\begin{frame}
\frametitle{公式}
$$M=m_{ubm}+Dy$$
\end{frame}

%%%%%%%%%%%%%%%%%%%%%%%%%%%%%%%%%%%%%%%%%%%%%%%%%%%%%%%%%%%%%%%%%%%%%%%%%%%%%%%

\section{方法三}

\begin{frame}
\frametitle{Block}
\begin{block}{Block1}
	\begin{itemize}
		\item 方法三
        \item 方法三
	\end{itemize}
\end{block}
\pause
\begin{block}{Block1}
	\begin{itemize}
		\item 方法三
        \item 方法三
	\end{itemize}
\end{block}
\end{frame}

%%%%%%%%%%%%%%%%%%%%%%%%%%%%%%%%%%%%%%%%%%%%%%%%%%%%%%%%%%%%%%%%%%%%%%%%%%%%%%%

\section{总结与展望}

\begin{frame}
\frametitle{实验结果}
\begin{columns}
\begin{column}{0.50\textwidth}
\epsfig{figure=ustc_logo_fig_new,width=\textwidth}
\end{column}
\begin{column}{0.50\textwidth}
\begin{block}{结论}
\begin{itemize}
\item 实验结果如左图所示
\item 实验结果如左图所示
\item 实验结果如左图所示
\item 实验结果如左图所示
\item 实验结果如左图所示
\item 实验结果如左图所示
\end{itemize}
\end{block}
\end{column}
\end{columns}
\end{frame}


%%%%%%%%%%%%%%%%%%%%%%%%%%%%%%%%%%%%%%%%%%%%%%%%%%%%%%%%%%%%%%%%%%%%%%%%%%%%%%%


\begin{frame}
\centerline{\Large 谢谢!}
\end{frame}

\end{document}
